% ===================================================================
%  پیوست الف — فهرست نمادها و اختصارها
%  هدف: یک‌جا آوردن همه‌ی نمادهای پرکاربرد تا خواننده (به‌ویژه در
%  ارائه‌ی کلاسی) بدون گشتن در متن، معنا و یکای هر کمیت را ببیند.
%  اعداد و تعریف‌ها با بخش‌های ۳ و ۴ هم‌خوان‌اند.
% ===================================================================
\appendix
\section*{پیوست الف: فهرست نمادها و اختصارها}
\label{sec:appendix-nomenclature}
\addcontentsline{toc}{section}{پیوست الف: فهرست نمادها و اختصارها}

جدول~\ref{tab:nomenclature} نمادها و اختصارهای پرکاربرد این مقاله را یک‌جا
گرد آورده است. سه کمیت نخست ($\Fp$، $T$ و $\eta_a$) سه ستون فقرات تحلیل
این پژوهش‌اند: نرخ واپاشی کل، کانال تابشی و نسبت میان آن دو.

\begin{table}[htbp]
  \centering
  \caption{فهرست نمادها و اختصارهای پرکاربرد.}
  \label{tab:nomenclature}
  \begin{tabular}{lll}
    \toprule
    نماد / اختصار & معنا & یکا / توضیح \\
    \midrule
    $\Fp$ & فاکتور پورسل: نسبت نرخ واپاشی کل به نرخ در فضای آزاد
          & بی‌بعد (رابطه~\eqref{eq:Fp}) \\
    $T$ & ضریب تقویت توان تابشی میدان دور
          & بی‌بعد (رابطه~\eqref{eq:T}) \\
    $\eta_a$ & بازده کوانتومی ظاهری، $\eta_a = T/\Fp$
          & بی‌بعد، گزارش‌شده به درصد \\
    $P_0$ & توان گسیل‌شده‌ی دوقطبی در فضای همگن & وات \\
    $P_{\text{rad}}$ & توان تابش‌شده به میدان دوردست & وات \\
    $P_{\text{loss}}$ & توان اتلافی غیرتابشی (اهمی) در فلز & وات \\
    $\Gamma_{\text{tot}}$، $\Gamma_0$ & نرخ واپاشی کل / نرخ در فضای همگن
          & $\text{s}^{-1}$ \\
    $\lambda_{\text{res}}$ & طول‌موج قله‌ی فاکتور پورسل & $\nm{607.9}$ در این کار \\
    $\lambda_{\text{rad}}$ & طول‌موج قله‌ی توان تابشی & $\nm{615.1}$ در این کار \\
    $L$، $D$، $R$ & طول کل، قطر و شعاع کلاهک نانومیله & $\nm{60}$، $\nm{20}$، $\nm{10}$ \\
    $d$ & فاصله‌ی دوقطبی تا راس کلاهک & $\nm{5}$ \\
    $l$، $m$ & مرتبه‌ی چندقطبی کروی و شماره‌ی آزیموتالی & بی‌بعد \\
    $\beta$ & ثابت انتشار محوری در هندسه‌ی استوانه‌ای & $\text{rad/m}$ \\
    \midrule
    LDOS & چگالی موضعی حالات فوتونی & --- \\
    LSPR & پلاسمون سطحی موضعی & --- \\
    FDTD & تفاضل محدود در حوزه زمان & --- \\
    PML & لایه‌ی جاذب کاملاً منطبق (مرز محاسباتی) & --- \\
    QNM & مد شبه‌نرمال (مدهای کاواک با فرکانس مختلط) & --- \\
    \bottomrule
  \end{tabular}
\end{table}
